\section{Variáveis Aleatórias Discretas}

\begin{def:variavel discreta}
Se $X$ é uma variável aleatória, então $X$ é discreta caso a sua FDA seu suporte
seja enumerável.
\end{def:variavel discreta}

Dessa definição, podemos determinar que existe um número contável de pontos
$x \in \mathbb{R}$ cuja toda vizinhança aberta $U$ centrada neles tenha a propriedade
de $P(X \in U) > 0$. Esse fato diz que podemos associar a cada ponto uma determinada
``massa'' de probabilidade. A função que realiza esse mapeamento é chamada de
\textbf{função de massa de probabilidade}.

\begin{def:fmp}
Seja $X: \Omega \to \mathbb{R}$ uma v.a. discreta. A função $p_X: \mathbb{R} \to \mathbb{R}$
é a função de massa de probabilidade de $X$ se, e somente se,
\[
	P(X = x) = p_X(x)
\]
\end{def:fmp}

\begin{prop:fmp suporte}
Se $X: \Omega \to \mathbb{R}$ é uma v.a. discreta com f.m.p. igual a $p_X$, então
\[
	\sum_{x \in \Omega_X} p_X(x) = 1
\]
\begin{proof}
	Como $\Omega_x$ é enumerável, então presuma que $\Omega_X = \{\alpha_1, \ldots, \alpha_n\}$
	e que $\Parts{\Omega} = \{E_1, \ldots, E_n\}$ seja uma partição de $\Omega$ tal que
	\[
		\forall i \in [n],\ (P(E_i) = p_X(\alpha_i))
	\]
	Com isso,
	\[
		P\left(\bigcup_{i=1}^n E_i\right) =
		\sum_{i = 1}^n P(E_i) =
		\sum_{i = 1}^n p_X(\alpha_i)
	\]
	donde se conclui
	\[
		\sum_{i = 1}^n p_X(\alpha_i) = 1
	\]
\end{proof}
\end{prop:fmp suporte}

Pela definição de variáveis discretas, então digamos que $\{x_1, \ldots, x_n\}$ seja uma
ordenação de $\Omega_X$ de tal modo
\[
	F(X) = \begin{cases}
		0,          & \text{ se } x < x_1              \\
		F(x_1),     & \text{ se } x_1 \leq x < x_2     \\
		\vdots                                         \\
		F(x_{n-1}), & \text{ se } x_{n-1} \leq x < x_n \\
		1,          & \text{ se } x_n \leq x           \\
	\end{cases}.
\]
A função $F$ é, pois, uma função constante por partes que sempre incluem a extremidade
esquerda de seu domínio devido à propriedade necessária de uma FDA sempre ser contínua
à direita. Uma vez que
\begin{gather*}
	P(X = x_i) = P(X \leq x_i) - P(X < x_i)\\
	P(X = x_i) = F(x_i) - \lim_{x \to x_i^-} F(x)\\
	P(X = x_i) = F(x_i) - F(x_{i-1}),
\end{gather*}
então $p_X(x_i)$ é o salto entre duas partes consecutivas da FDA de $X$.

Exemplos de variáveis aleatórias discretas incluem
\begin{itemize}
	\item Número de caras obtidas em 10 lançamentos.
	\item Soma dos lados obtidos no lançamento de dados.
	\item Número de clientes em uma loja em um dia.
\end{itemize}
No primeiro exemplo, as realizações possíveis são os inteiros de 0 a 10. Nos dois seguintes,
são todos os naturais. Todos esses são conjunto enumeráveis, cujos pontos isoladamente
concentram massa de probabilidade, isto é, um valor efetivo de probabilidade associado à
ao evento da variável de ter realização naquele ponto.
